\begin{table*}[t]
\centering
\caption{Mainline comparison and ablation design. A4--A9 use the same network;
checkmarks indicate active loss groups. A7--A9 add trunk-twist objectives on top
of A6 (per-view-peak ROM anchor, an optional rotation-parameterised twist
residual, and a per-view observed twist-rate anchor).}
\label{tab:ablation-design}
\scriptsize
\setlength{\tabcolsep}{3pt}
\begin{tabularx}{\linewidth}{c X c c c c c}
\toprule
ID & Method & \makecell{Learned\\residual} &
\makecell{Recovery, identity,\\and structure} &
\makecell{Rotation and\\temporal} & \makecell{Complete-cycle\\ROM} &
\makecell{Trunk\\twist} \\
\midrule
A0 & Face stream only & -- & -- & -- & -- & -- \\
A1 & Side stream in face reference & -- & -- & -- & -- & -- \\
A2 & Canonical arithmetic mean & -- & -- & -- & -- & -- \\
A3 & Deterministic quality-weighted mean & -- & -- & -- & -- & -- \\
A4 & Learned spatial-objective fusion & \checkmark & \checkmark & -- & -- & -- \\
A5 & Learned rotation-temporal fusion & \checkmark & \checkmark & \checkmark & -- & -- \\
A6 & Full rotation-aware self-supervised fusion & \checkmark & \checkmark & \checkmark & \checkmark & -- \\
A7 & A6 + per-view-peak ROM anchor & \checkmark & \checkmark & \checkmark & \checkmark & \checkmark \\
A8 & A7 + rotation-parameterised twist residual & \checkmark & \checkmark & \checkmark & \checkmark & \checkmark \\
A9 & A8 + observed twist-rate anchor & \checkmark & \checkmark & \checkmark & \checkmark & \checkmark \\
\bottomrule
\end{tabularx}
\end{table*}
